
\documentclass{article}
\usepackage{colortbl}
\usepackage{makecell}
\usepackage{multirow}
\usepackage{supertabular}

\begin{document}

\newcounter{utterance}

\centering \large Interaction Transcript for game `referencegame', experiment `random\_grids', episode 8 with reflexion{-}Qwen3.5{-}27B.
\vspace{24pt}

{ \footnotesize  \setcounter{utterance}{1}
\setlength{\tabcolsep}{0pt}
\begin{supertabular}{c@{$\;$}|p{.15\linewidth}@{}p{.15\linewidth}p{.15\linewidth}p{.15\linewidth}p{.15\linewidth}p{.15\linewidth}}
    \# & $\;$A & \multicolumn{4}{c}{Game Master} & $\;\:$B\\
    \hline

    \theutterance \stepcounter{utterance}  
    & & \multicolumn{4}{p{0.6\linewidth}}{
        \cellcolor[rgb]{0.9,0.9,0.9}{
            \makecell[{{p{\linewidth}}}]{
                \texttt{\tiny{[P1$\langle$GM]}}
                \texttt{You have attempted this same task before and did not succeed. These are the plans you wrote after those attempts.} \\
\texttt{Trial \#1: New plan: My previous description failed because the features I selected (top{-}left empty, middle filled cells in top/bottom rows) were not unique to the target grid; Distractor 2 also shared these exact characteristics. To succeed, I must identify a feature present in the target grid but absent in *both* distractors. Specifically, I will analyze the second row: the target has a filled cell at the far left (column 1) and far right (column 5) with empty cells in between, whereas Distractor 1 has an empty left cell and Distractor 2 has a filled left cell but an empty right cell. My new strategy is to construct the expression based on this unique configuration in the second row (e.g., "The grid where the second row starts and ends with a filled cell") to ensure strict differentiation.} \\
\\ 
\\ 
\texttt{You are given three grids, where each of them is 5 by 5 in size.} \\
\texttt{Grids have empty cells marked with "▢" and filled cells marked with "X".} \\
\texttt{Your task is to generate a referring expression that best describes the target grid while distinguishing it from the two other distractor grids.} \\
\texttt{The first grid is the target grid, and the following two grids are the distractors.} \\
\\ 
\texttt{Target grid:} \\
\\ 
\texttt{□ □ X X □} \\
\texttt{X □ □ X X} \\
\texttt{□ X X □ X} \\
\texttt{X X □ □ X} \\
\texttt{□ X X □ □} \\
\\ 
\texttt{Distractor grid 1:} \\
\\ 
\texttt{X □ X X □} \\
\texttt{□ X □ X X} \\
\texttt{X X □ □ X} \\
\texttt{□ X X □ X} \\
\texttt{X □ X □ □} \\
\\ 
\texttt{Distractor grid 2:} \\
\\ 
\texttt{□ X X □ X} \\
\texttt{X □ X X □} \\
\texttt{□ X □ X X} \\
\texttt{X □ X □ X} \\
\texttt{□ X X □ □} \\
\\ 
\texttt{Instruction: Describe the target grid.} \\
\texttt{Generate the referring expression starting with the tag "Expression: " for the given target grid. Omit any other text.} \\
            }
        }
    }
    & \\ \\

    \theutterance \stepcounter{utterance}  
    & \multicolumn{4}{p{0.6\linewidth}}{
        \cellcolor[rgb]{0.9,0.9,0.9}{
            \makecell[{{p{\linewidth}}}]{
                \texttt{\tiny{[P1$\rangle$GM]}}
                \texttt{Expression: The grid where the second row starts and ends with a filled cell.} \\
            }
        }
    }
    & & \\ \\

    \theutterance \stepcounter{utterance}  
    & & & \multicolumn{2}{p{0.3\linewidth}}{
        \cellcolor[rgb]{0.9,0.9,0.9}{
            \makecell[{{p{\linewidth}}}]{
                \texttt{\tiny{[GM$|$GM]}}
                \texttt{Expression: The grid where the second row starts and ends with a filled cell.} \\
            }
        }
    }
    & & \\ \\

    \theutterance \stepcounter{utterance}  
    & & \multicolumn{4}{p{0.6\linewidth}}{
        \cellcolor[rgb]{0.9,0.9,0.9}{
            \makecell[{{p{\linewidth}}}]{
                \texttt{\tiny{[GM$\rangle$P2]}}
                \texttt{You have attempted this same task before and did not succeed. These are the plans you wrote after those attempts.} \\
\texttt{Trial \#1: New plan: I incorrectly selected the first grid because I failed to verify all conditions of the referring expression against the grid data. Specifically, while the first grid has an empty top{-}left cell and two consecutive filled cells in the middle of the top row, its bottom row (□ X X □ □) has the consecutive filled cells in the middle, which actually matches the description. However, I need to re{-}evaluate the "middle" constraint more strictly or check the other grids to ensure no better fit exists. Upon closer inspection of the third grid: Top{-}left is empty (□), top row has "X X" in the middle (columns 3 and 4), and bottom row has "X X" in the middle (columns 2 and 3). Wait, let's re{-}read the grids carefully.} \\
\texttt{Grid 1 Top: □ X X □ X (Consecutive X at 2,3). Bottom: □ X X □ □ (Consecutive X at 2,3).} \\
\texttt{Grid 3 Top: □ □ X X □ (Consecutive X at 3,4). Bottom: □ X X □ □ (Consecutive X at 2,3).} \\
\texttt{The expression says "two consecutive filled cells in the middle". In a 5{-}cell row, "middle" usually implies the central position or the most central pair.} \\
\texttt{Let's re{-}evaluate the specific failure: I likely assumed the first grid was correct without rigorously checking the "middle" definition against the third grid. The third grid's top row has the pair at indices 3 and 4, which is slightly right of center, whereas the first grid has them at 2 and 3.} \\
\texttt{Actually, looking at the third grid again: Top row is `□ □ X X □`. The consecutive cells are at positions 3 and 4. Bottom row is `□ X X □ □`. The consecutive cells are at positions 2 and 3.} \\
\texttt{The first grid: Top row `□ X X □ X` (positions 2,3). Bottom row `□ X X □ □` (positions 2,3).} \\
\texttt{The expression requires: 1. Top{-}left empty. 2. Top row has two consecutive filled cells in the middle. 3. Bottom row has two consecutive filled cells in the middle.} \\
\texttt{If "middle" means the exact center of the 5 cells, there is no single center cell for a pair, so it likely means the pair is centered (positions 2{-}3 or 3{-}4).} \\
\texttt{My mistake was likely a hasty visual match on the first grid without confirming if the third grid was a better fit or if I misread the third grid's bottom row.} \\
\texttt{Correct Strategy: I will explicitly parse the grid row by row.} \\
\texttt{1. Check Top{-}Left: Must be "□". (Grid 1: Yes, Grid 2: No, Grid 3: Yes} \\
\\ 
\\ 
\texttt{You are given three grids, where each of them is 5 by 5 in size.} \\
\texttt{Grids have empty cells marked with "▢" and filled cells marked with "X".} \\
\texttt{You are also given a referring expression that describes one of the given grids.} \\
\texttt{Your task is to select the grid that matches the given referring expression.} \\
\texttt{Generate only the number (in text) of the grid that matches the given expression by selecting first, second, or third.} \\
\\ 
\texttt{First:} \\
\\ 
\texttt{□ X X □ X} \\
\texttt{X □ X X □} \\
\texttt{□ X □ X X} \\
\texttt{X □ X □ X} \\
\texttt{□ X X □ □} \\
\\ 
\texttt{Second:} \\
\\ 
\texttt{X □ X X □} \\
\texttt{□ X □ X X} \\
\texttt{X X □ □ X} \\
\texttt{□ X X □ X} \\
\texttt{X □ X □ □} \\
\\ 
\texttt{Third:} \\
\\ 
\texttt{□ □ X X □} \\
\texttt{X □ □ X X} \\
\texttt{□ X X □ X} \\
\texttt{X X □ □ X} \\
\texttt{□ X X □ □} \\
\\ 
\texttt{Expression: The grid where the second row starts and ends with a filled cell.} \\
\texttt{Question: Which grid does the expression refer to?} \\
\texttt{Start with the tag "Answer: ", followed by your selection. Omit any other text.} \\
            }
        }
    }
    & \\ \\

    \theutterance \stepcounter{utterance}  
    & & & \multicolumn{4}{p{0.6\linewidth}}{
        \cellcolor[rgb]{0.9,0.9,0.9}{
            \makecell[{{p{\linewidth}}}]{
                \texttt{\tiny{[GM$\langle$P2]}}
                \texttt{Answer: second} \\
            }
        }
    }
     \\ \\

    \theutterance \stepcounter{utterance}  
    & & & \multicolumn{2}{p{0.3\linewidth}}{
        \cellcolor[rgb]{0.9,0.9,0.9}{
            \makecell[{{p{\linewidth}}}]{
                \texttt{\tiny{[GM$|$GM]}}
                \texttt{second} \\
            }
        }
    }
    & & \\ \\

    \theutterance \stepcounter{utterance}  
    & \multicolumn{4}{p{0.6\linewidth}}{
        \cellcolor[rgb]{0.95,0.95,0.95}{
            \makecell[{{p{\linewidth}}}]{
                \texttt{\tiny{[P2$\rightarrow$P1]}}
                \texttt{Post{-}episode Reflexion (not sent to GM):} \\
\texttt{New plan: My previous description failed because the feature I selected (second row starting and ending with a filled cell) was not unique; Distractor 1 also has a filled cell at the start and end of its second row. To succeed, I must identify a feature present in the target grid but absent in *both* distractors. Specifically, I will analyze the first row: the target grid has empty cells in the first two columns (□ □), whereas Distractor 1 starts with a filled cell (X) and Distractor 2 has a filled cell in the second column (□ X). My new strategy is to construct the expression based on this unique configuration in the first row (e.g., "The grid where the first two cells of the first row are empty") to ensure strict differentiation.} \\
            }
        }
    }
    & & \\ \\

    \theutterance \stepcounter{utterance}  
    & & & \multicolumn{4}{p{0.6\linewidth}}{
        \cellcolor[rgb]{0.95,0.95,0.95}{
            \makecell[{{p{\linewidth}}}]{
                \texttt{\tiny{[P1$\rightarrow$P2]}}
                \texttt{Post{-}episode Reflexion (not sent to GM):} \\
\texttt{New plan: I incorrectly selected the second grid because I failed to verify the specific condition "second row starts and ends with a filled cell" against the actual grid data. In the second grid, the second row is `□ X □ X X`, which starts with an empty cell (□), not a filled one (X). I likely misread the row index or the cell values. Next time, I will explicitly parse the second row of each grid:} \\
\texttt{1. Check Grid 1 Row 2: `X □ X X □` (Starts with X, ends with □) {-}> Fail.} \\
\texttt{2. Check Grid 2 Row 2: `□ X □ X X` (Starts with □, ends with X) {-}> Fail.} \\
\texttt{3. Check Grid 3 Row 2: `X □ □ X X` (Starts with X, ends with X) {-}> Pass.} \\
\texttt{I will ensure I read the exact characters at the first and last positions of the specified row before making a selection.} \\
            }
        }
    }
     \\ \\

\end{supertabular}
}

\end{document}
